\documentclass[11pt,a4paper]{article}
\usepackage[margin=1in]{geometry}
\usepackage{amsmath,amssymb,amsthm}
\usepackage[T1]{fontenc}
\usepackage[utf8]{inputenc}
\usepackage{mathptmx}
\usepackage{hyperref}
\usepackage{booktabs}
\usepackage{longtable}
\usepackage{graphicx}
\usepackage{enumitem}
\usepackage{xcolor}
\usepackage{array}
\usepackage{microtype}
\usepackage{parskip}
\usepackage{titlesec}
\usepackage{textcomp}
\usepackage{newunicodechar}

% Unicode character definitions
\DeclareUnicodeCharacter{2194}{\ensuremath{\leftrightarrow}}
\DeclareUnicodeCharacter{2192}{\ensuremath{\to}}
\DeclareUnicodeCharacter{2190}{\ensuremath{\leftarrow}}
\DeclareUnicodeCharacter{2248}{\ensuremath{\approx}}
\DeclareUnicodeCharacter{2264}{\ensuremath{\leq}}
\DeclareUnicodeCharacter{2265}{\ensuremath{\geq}}
\DeclareUnicodeCharacter{2261}{\ensuremath{\equiv}}
\DeclareUnicodeCharacter{2260}{\ensuremath{\neq}}
\DeclareUnicodeCharacter{2208}{\ensuremath{\in}}
\DeclareUnicodeCharacter{2205}{\ensuremath{\emptyset}}
\DeclareUnicodeCharacter{221E}{\ensuremath{\infty}}
\DeclareUnicodeCharacter{2297}{\ensuremath{\otimes}}
\DeclareUnicodeCharacter{2282}{\ensuremath{\subset}}
\DeclareUnicodeCharacter{2283}{\ensuremath{\supset}}
\DeclareUnicodeCharacter{00D7}{\ensuremath{\times}}
\DeclareUnicodeCharacter{00B1}{\ensuremath{\pm}}
\DeclareUnicodeCharacter{00B7}{\ensuremath{\cdot}}
\DeclareUnicodeCharacter{2124}{\ensuremath{\mathbb{Z}}}
\DeclareUnicodeCharacter{2212}{\ensuremath{-}}
\DeclareUnicodeCharacter{223C}{\ensuremath{\sim}}
\DeclareUnicodeCharacter{210F}{\ensuremath{\hbar}}
\DeclareUnicodeCharacter{2713}{\checkmark}
\DeclareUnicodeCharacter{2717}{\ensuremath{\times}}
\DeclareUnicodeCharacter{00B2}{\ensuremath{^2}}
\DeclareUnicodeCharacter{00B3}{\ensuremath{^3}}
\DeclareUnicodeCharacter{00B9}{\ensuremath{^1}}
\DeclareUnicodeCharacter{2076}{\ensuremath{^6}}
\DeclareUnicodeCharacter{207B}{\ensuremath{^-}}
\DeclareUnicodeCharacter{2077}{\ensuremath{^7}}
\DeclareUnicodeCharacter{00B0}{\ensuremath{^\circ}}

% Greek uppercase
\DeclareUnicodeCharacter{0393}{\ensuremath{\Gamma}}
\DeclareUnicodeCharacter{0394}{\ensuremath{\Delta}}
\DeclareUnicodeCharacter{039B}{\ensuremath{\Lambda}}
\DeclareUnicodeCharacter{03A3}{\ensuremath{\Sigma}}

% Greek lowercase
\DeclareUnicodeCharacter{03B1}{\ensuremath{\alpha}}
\DeclareUnicodeCharacter{03B2}{\ensuremath{\beta}}
\DeclareUnicodeCharacter{03B3}{\ensuremath{\gamma}}
\DeclareUnicodeCharacter{03B4}{\ensuremath{\delta}}
\DeclareUnicodeCharacter{03B5}{\ensuremath{\varepsilon}}
\DeclareUnicodeCharacter{03B7}{\ensuremath{\eta}}
\DeclareUnicodeCharacter{03B8}{\ensuremath{\theta}}
\DeclareUnicodeCharacter{03BA}{\ensuremath{\kappa}}
\DeclareUnicodeCharacter{03BB}{\ensuremath{\lambda}}
\DeclareUnicodeCharacter{03BC}{\ensuremath{\mu}}
\DeclareUnicodeCharacter{03BD}{\ensuremath{\nu}}
\DeclareUnicodeCharacter{03C0}{\ensuremath{\pi}}
\DeclareUnicodeCharacter{03C1}{\ensuremath{\rho}}
\DeclareUnicodeCharacter{03C3}{\ensuremath{\sigma}}
\DeclareUnicodeCharacter{03C4}{\ensuremath{\tau}}
\DeclareUnicodeCharacter{03C6}{\ensuremath{\varphi}}
\DeclareUnicodeCharacter{03C7}{\ensuremath{\chi}}
\DeclareUnicodeCharacter{03C8}{\ensuremath{\psi}}
\DeclareUnicodeCharacter{03C9}{\ensuremath{\omega}}

% Superscripts and subscripts
\DeclareUnicodeCharacter{2070}{\ensuremath{^0}}
\DeclareUnicodeCharacter{2075}{\ensuremath{^5}}
\DeclareUnicodeCharacter{2078}{\ensuremath{^8}}
\DeclareUnicodeCharacter{2074}{\ensuremath{^4}}
\DeclareUnicodeCharacter{207A}{\ensuremath{^+}}
\DeclareUnicodeCharacter{2080}{\ensuremath{_0}}
\DeclareUnicodeCharacter{2081}{\ensuremath{_1}}
\DeclareUnicodeCharacter{2082}{\ensuremath{_2}}
\DeclareUnicodeCharacter{2083}{\ensuremath{_3}}
\DeclareUnicodeCharacter{2084}{\ensuremath{_4}}

% Math symbols
\DeclareUnicodeCharacter{2202}{\ensuremath{\partial}}
\DeclareUnicodeCharacter{221A}{\ensuremath{\sqrt{}}}
\DeclareUnicodeCharacter{226A}{\ensuremath{\ll}}
\DeclareUnicodeCharacter{226B}{\ensuremath{\gg}}
\DeclareUnicodeCharacter{21D2}{\ensuremath{\Rightarrow}}
\DeclareUnicodeCharacter{27F9}{\ensuremath{\Longrightarrow}}
\DeclareUnicodeCharacter{27FA}{\ensuremath{\Longleftrightarrow}}

% Modifier letters
\DeclareUnicodeCharacter{02B2}{\ensuremath{^j}}
\DeclareUnicodeCharacter{0304}{\={}}

% Additional symbols that may appear
\DeclareUnicodeCharacter{2113}{\ensuremath{\ell}}
\DeclareUnicodeCharacter{2229}{\ensuremath{\cap}}
\DeclareUnicodeCharacter{222A}{\ensuremath{\cup}}
\DeclareUnicodeCharacter{2227}{\ensuremath{\wedge}}
\DeclareUnicodeCharacter{2228}{\ensuremath{\vee}}
\DeclareUnicodeCharacter{00AC}{\ensuremath{\neg}}

% hyperref setup
\hypersetup{
  colorlinks=true,
  linkcolor=blue!60!black,
  citecolor=blue!60!black,
  urlcolor=blue!60!black,
  pdftitle={The Structural Preconditions for Complexity: Why the Laws of Physics Guarantee the Possibility of Organic Chemistry},
  pdfauthor={Alex Maybaum}
}

% Section formatting - no auto-numbering (source has its own)
\setcounter{secnumdepth}{0}
\titleformat{\section}{\Large\bfseries}{}{0em}{}
\titleformat{\subsection}{\large\bfseries}{}{0em}{}
\titleformat{\subsubsection}{\normalsize\bfseries}{}{0em}{}

% Tight list spacing
\providecommand{\tightlist}{\setlength{\itemsep}{0pt}\setlength{\parskip}{0pt}}

% Fix for pandoc's CSLReferences
\newlength{\cslhangindent}
\setlength{\cslhangindent}{1.5em}
\newenvironment{CSLReferences}[2]{}{}

\title{The Structural Preconditions for Complexity: Why the Laws of Physics Guarantee the Possibility of Organic Chemistry}
\author{Alex Maybaum}
\date{March 2026}

\begin{document}

\hypertarget{the-structural-preconditions-for-complexity}{%
\section{The Structural Preconditions for Complexity}\label{the-structural-preconditions-for-complexity}}

\hypertarget{why-the-laws-of-physics-guarantee-the-possibility-of-organic-chemistry}{%
\subsubsection{Why the Laws of Physics Guarantee the Possibility of Organic Chemistry}\label{why-the-laws-of-physics-guarantee-the-possibility-of-organic-chemistry}}

\textbf{Author:} Alex Maybaum\\
\textbf{Date:} March 2026\\
\textbf{Status:} DRAFT PRE-PRINT\\
\textbf{Classification:} Theoretical Physics / Foundations / Astrobiology

\begin{center}\rule{0.5\linewidth}{0.5pt}\end{center}

\hypertarget{abstract}{%
\subsection{Abstract}\label{abstract}}

The Observational Incompleteness (OI) framework derives the Standard Model's structure --- gauge group SU(3) × SU(2) × U(1), three fermion generations, one Higgs doublet, chiral weak interactions, and θ̄ = 0 --- from four axioms about embedded observation {[}1, 2{]}. We trace the consequences of this derived structure through atomic physics, chemistry, and prebiotic selection, establishing a structural chain from (S, φ) to the preconditions for organic chemistry. The chain is entirely parameter-free: the spatial dimension d = 3 determines the orbital algebra SO(3), which determines the periodic table's architecture, which determines carbon's tetrahedral bonding geometry. Three generations guarantee CP violation and hence baryogenesis. The Higgs mechanism guarantees a mass hierarchy with light fermions. Chiral SU(2) produces a universal parity-violating energy difference between mirror-image molecules, with the sign fixed by the partition structure.

The structural chain establishes that the \emph{possibility} of carbon-based organic chemistry, matter-antimatter asymmetry, chiral molecular selection, a scale hierarchy separating nuclear and atomic physics, water with anomalous solvent properties, and a thermal window for dynamic chemistry is a theorem about embedded observation --- not a contingent fact about our universe. The \emph{actuality} of chemistry depends on 18 parameter values within the derived structure. Using existing constraints from nuclear physics, atomic physics, and chemistry, we estimate the viable fraction of the parameter space at \(\gtrsim 10\%\) for the most constrained parameters (the fine-structure constant and the quark mass ratio), substantially larger than fine-tuning arguments assume when both structure and parameters are varied. If the viable fraction is confirmed to be large, the fine-tuning problem dissolves: the laws are derived, the preconditions are structural, and the parameters are not special. Furthermore, the combinatorial diversity that carbon chemistry in d = 3 necessarily produces (\(\sim 20^N\) possible peptides of length \(N\)) exceeds both the autocatalytic threshold and von Neumann's self-reproducing automaton threshold by tens of orders of magnitude, making self-sustaining, self-reproducing chemical networks a statistical expectation rather than a contingent accident. Template replication follows from the structural intersection of linear polymers, complementary pairing, and catalytic activity --- all independently guaranteed by the orbital algebra. Darwinian evolution is the inevitable dynamics of imperfect template replication. The chain extends through information processing (generically fitness-enhancing), neural computation (structurally possible), and artificial intelligence (structurally guaranteed given general intelligence). The one contingent step is general intelligence itself; every other link is structural, statistical, or inevitable given the preceding link. The characterization theorem imposes a permanent structural floor on all embedded observers --- biological or artificial --- while leaving no ceiling on their complexity.

\begin{center}\rule{0.5\linewidth}{0.5pt}\end{center}

\hypertarget{introduction}{%
\subsection{1. Introduction}\label{introduction}}

The fine-tuning problem in physics rests on a counterfactual: if the laws of physics or the values of fundamental constants were different, the universe would not permit complex structures, and we would not exist to observe it. The anthropic principle elevates this observation into an explanatory principle --- our existence selects the laws from an ensemble of possibilities.

The OI framework {[}1, 2{]} challenges the premise. The companion papers prove that an embedded observer in a deterministic system with conditions C1--C3 (coupling, slow bath, sufficient capacity) necessarily describes the visible sector using quantum mechanics {[}1{]}, and that the dynamics uniquely selected by this requirement determines the Standard Model's complete structure {[}2{]}. If the laws are derived rather than contingent, the largest source of anthropic variation --- the structure of the laws themselves --- is eliminated.

This paper traces the consequences. §2 establishes the structural chain from (S, φ) to the full preconditions for organic chemistry --- orbital structure, the periodic table, carbon's bonding geometry, the nuclear-atomic scale hierarchy, water's solvent properties, and the thermal window for dynamic chemistry. §3 estimates the viable fraction of the parameter space using existing literature. §4 derives the chirality chain from the partition to biological homochirality. §5 extends the chain through autocatalytic networks, template replication, Darwinian evolution, information processing, and artificial intelligence, identifying the one contingent step and the structural limits on all embedded observers.

\begin{center}\rule{0.5\linewidth}{0.5pt}\end{center}

\hypertarget{the-structural-chain}{%
\subsection{2. The Structural Chain}\label{the-structural-chain}}

\hypertarget{from-s-ux3c6-to-orbitals}{%
\subsubsection{2.1 From (S, φ) to orbitals}\label{from-s-ux3c6-to-orbitals}}

The framework derives d = 3 {[}2, §6{]} from four independent self-consistency filters (propagating gravity, stable matter, the dark sector concordance, and renormalizability). In three spatial dimensions, the rotation group is SO(3), whose irreducible representations are labeled by angular momentum quantum number \(l = 0, 1, 2, 3, \ldots\) with \((2l+1)\) states each. These are the atomic orbitals: s (\(l = 0\), spherical, 1 orientation), p (\(l = 1\), dumbbell, 3 orientations), d (\(l = 2\), 5 orientations), f (\(l = 3\), 7 orientations).

This is the \emph{only} orbital algebra consistent with embedded observation. In \(d = 2\), angular momentum has a single component --- only circular harmonics exist, with no three-dimensional bonding geometry. In \(d \geq 4\), the Coulomb potential \(V(r) \sim r^{2-d}\) does not support stable bound states {[}3, 4{]}: classical orbits spiral into the nucleus, and the quantum Hamiltonian is unbounded below.

\hypertarget{from-orbitals-to-the-periodic-table}{%
\subsubsection{2.2 From orbitals to the periodic table}\label{from-orbitals-to-the-periodic-table}}

The framework derives spin-1/2 fermions from the staggered structure on the d = 3 lattice {[}2, §7.2{]}. The spin-statistics theorem --- a consequence of the emergent QFT's Lorentz invariance {[}2, §5{]} --- gives the Pauli exclusion principle: each orbital holds at most 2 electrons (spin up and spin down). The shell capacities follow:

\[s: 2, \quad p: 6, \quad d: 10, \quad f: 14\]

The periodic table's period lengths --- 2, 8, 8, 18, 18, 32 --- are determined by the filling order of these shells. This architecture is entirely structural: it depends on d = 3, spin-1/2, and the Pauli principle, all of which are derived.

\hypertarget{carbons-bonding-geometry}{%
\subsubsection{2.3 Carbon's bonding geometry}\label{carbons-bonding-geometry}}

Element 6 (carbon) has electron configuration 1s² 2s² 2p² --- four valence electrons. In d = 3, four valence electrons admit three hybridization modes:

\emph{sp³ hybridization:} Four equivalent orbitals pointing to the vertices of a regular tetrahedron (bond angle 109.5°). This is the geometry of methane (CH₄), diamond, and the backbone of all organic macromolecules --- proteins, nucleic acids, lipids, carbohydrates.

\emph{sp² hybridization:} Three planar orbitals at 120° plus one unhybridized p orbital forming a π bond. This is the geometry of graphite, benzene, and all aromatic chemistry --- the basis of most pharmaceutical and biological signaling molecules.

\emph{sp hybridization:} Two linear orbitals at 180° plus two π bonds. This is the geometry of acetylene and the basis of alkyne chemistry.

The \emph{possibility} of these hybridizations --- and hence the possibility of three-dimensional macromolecular architecture --- is a structural consequence of the angular momentum algebra in d = 3. No parameter values are involved.

\hypertarget{beyond-carbon-the-extended-structural-chain}{%
\subsubsection{2.4 Beyond carbon: the extended structural chain}\label{beyond-carbon-the-extended-structural-chain}}

Four additional structural results extend the chain beyond orbital chemistry.

\textbf{Matter-antimatter asymmetry.} Three generations {[}2, §7.7{]} give the CKM matrix with a physical CP-violating phase --- one of the three Sakharov conditions for baryogenesis {[}5{]}. The partition breaks P {[}2, §7.8, Theorem 13{]}. Combined with CP violation, baryogenesis is structurally possible. Without three generations, the CKM matrix has no CP-violating phase and the standard electroweak baryogenesis mechanism does not operate.

\textbf{Mass hierarchy.} The Higgs mechanism {[}2, §7.3{]} gives fermion masses through Yukawa couplings to a single doublet. The taste-breaking mechanism {[}2, §7.7{]} produces generation-dependent couplings generically, ensuring different generations have different masses. This guarantees the existence of light fermions (electron, u and d quarks) --- not as a coincidence, but as a structural feature of the staggered lattice.

\textbf{Nuclear binding.} SU(3) color confinement {[}2, §7.6{]} produces bound hadrons. The existence of multi-nucleon bound states (nuclei) requires specific parameter values (§3), but the \emph{mechanism} --- color confinement binding quarks into hadrons, residual strong force binding hadrons into nuclei --- is structural.

\textbf{Chiral chemistry.} Chiral SU(2) {[}2, §7.8{]} produces electroweak parity violation in all weak processes. This propagates to molecular physics as a parity-violating energy difference between mirror-image molecules (§4).

\hypertarget{the-scale-hierarchy}{%
\subsubsection{2.5 The scale hierarchy}\label{the-scale-hierarchy}}

The framework derives two independent gauge groups --- SU(3) and U(1) --- with independent coupling strengths corresponding to independent eigenvalues of the coupling matrix M {[}2, §7.6{]}. This guarantees two widely separated energy scales:

\[E_{\text{nuclear}} \sim \Lambda_{\text{QCD}} \sim \text{MeV}, \qquad E_{\text{atomic}} \sim \alpha^2 m_e c^2 \sim \text{eV}\]

The ratio \(E_{\text{nuclear}} / E_{\text{atomic}} \sim 10^6\) is partly parametric (it depends on the specific couplings), but the \emph{existence} of two independent scales is structural --- it follows from having two independent gauge groups. This separation has a profound consequence: nuclei are stable against chemical reactions. Chemistry rearranges electrons (eV) without touching nuclei (MeV). Atoms are permanent building blocks that can be assembled, disassembled, and reassembled without being destroyed.

A universe with a single gauge group would have one energy scale. Chemistry and nuclear physics would be inseparable. Every chemical reaction would risk nuclear transmutation. No stable building blocks, no reusable molecular components, no possibility of the combinatorial complexity that biology requires.

\hypertarget{water}{%
\subsubsection{2.6 Water}\label{water}}

Element 8 (oxygen): configuration 1s² 2s² 2p⁴ --- six valence electrons, two unpaired. The sp³-like hybridization in d = 3, distorted by two lone pairs, gives a bond angle of \textasciitilde104.5°. This is the geometry of water (H₂O).

Water's anomalous properties follow from this geometry. The bent structure creates a permanent dipole moment. The lone pairs enable hydrogen bonding --- a secondary interaction (\textasciitilde0.2 eV) intermediate between covalent bonds (\textasciitilde3 eV) and thermal energy (\textasciitilde0.025 eV at 300K). Hydrogen bonding produces: high specific heat (thermal buffering for chemical reactions), density maximum near 4°C (ice floats, insulating liquid water below), and exceptional solvent capability (dissolving ionic and polar molecules for solution-phase chemistry).

The \emph{possibility} of a molecule with these properties is structural: element 8's electron configuration in d = 3 with the derived orbital algebra guarantees a bent dihydride with lone pairs capable of hydrogen bonding. The \emph{existence} of a liquid phase at specific temperatures is parametric.

\hypertarget{the-thermal-window}{%
\subsubsection{2.7 The thermal window}\label{the-thermal-window}}

For chemistry to support complexity, a thermal window must exist where bonds are stable but reactions are dynamic. The structural chain guarantees this window is non-empty:

\emph{Stability:} Bond energies (\textasciitilde1--5 eV for covalent bonds in the derived orbital algebra) are set by the atomic energy scale \(\alpha^2 m_e c^2\).

\emph{Dynamics:} Thermal fluctuations at temperature \(T\) have energy \(kT\). For \(kT \ll E_{\text{bond}}\), bonds are stable; for \(kT > E_{\text{activation}}\) (where \(E_{\text{activation}}\) is the barrier height for a specific reaction), reactions proceed at finite rate.

\emph{Solvent:} A liquid phase requires a temperature between the freezing and boiling points of the solvent.

The scale hierarchy (§2.5) guarantees that \(E_{\text{bond}} / E_{\text{thermal}} \sim \alpha^2 m_e c^2 / kT \gg 1\) at any temperature where a liquid exists --- because the liquid phase occurs at temperatures far below bond dissociation. This means chemistry is inherently \emph{stable but dynamic} in any universe with two independent gauge groups and the derived orbital structure. The specific width of the thermal window depends on parameters, but its \emph{existence} is structural.

\hypertarget{summary-of-the-structural-chain}{%
\subsubsection{2.8 Summary of the structural chain}\label{summary-of-the-structural-chain}}

\[\text{(S, } \varphi\text{)} \xrightarrow{\text{derived}} d = 3 \xrightarrow{\text{SO(3)}} \text{orbitals} \xrightarrow{\text{Pauli}} \text{periodic table} \xrightarrow{\text{element 6}} \text{tetrahedral carbon}\]

\[\text{(S, } \varphi\text{)} \xrightarrow{\text{derived}} \text{3 generations} \xrightarrow{\text{CKM}} \text{CP violation} \xrightarrow{\text{Sakharov}} \text{baryogenesis possible}\]

\[\text{(S, } \varphi\text{)} \xrightarrow{\text{derived}} \text{Higgs} \xrightarrow{\text{taste-breaking}} \text{mass hierarchy} \xrightarrow{} \text{light fermions} \xrightarrow{} \text{stable atoms}\]

\[\text{(S, } \varphi\text{)} \xrightarrow{\text{derived}} \text{SU(3)} \xrightarrow{\text{confinement}} \text{hadrons} \xrightarrow{\text{residual force}} \text{nuclei possible}\]

\[\text{(S, } \varphi\text{)} \xrightarrow{\text{derived}} \text{SU(3) + U(1) independent} \xrightarrow{} \text{scale hierarchy} \xrightarrow{} \text{atoms are permanent building blocks}\]

\[\text{(S, } \varphi\text{)} \xrightarrow{\text{derived}} \text{element 8 in d = 3} \xrightarrow{} \text{bent H}_2\text{O} \xrightarrow{\text{H-bonding}} \text{anomalous solvent}\]

\[\text{(S, } \varphi\text{)} \xrightarrow{\text{derived}} \text{scale hierarchy + orbital structure} \xrightarrow{} \text{thermal window exists}\]

\[\text{(S, } \varphi\text{)} \xrightarrow{\text{derived}} \text{chiral SU(2)} \xrightarrow{\text{PVED}} \text{L-amino acid preference}\]

Every arrow is either a theorem from {[}1, 2{]} or a standard textbook result applied to the derived structure. No parameter values enter. The possibility of a universe with atoms, nuclei, a periodic table, organic chemistry, matter-antimatter asymmetry, and chirally selected molecules is a theorem about embedded observation.

\begin{center}\rule{0.5\linewidth}{0.5pt}\end{center}

\hypertarget{the-viable-parameter-space}{%
\subsection{3. The Viable Parameter Space}\label{the-viable-parameter-space}}

\hypertarget{the-question}{%
\subsubsection{3.1 The question}\label{the-question}}

The structural chain establishes that organic chemistry is \emph{possible}. For it to be \emph{actual}, the 18 free parameters of the SM must fall in a viable region. How large is this region?

Traditional fine-tuning studies {[}6, 7, 8{]} vary both the structure and the parameters. The OI framework fixes the structure, reducing the question to parameters alone. This sharpens the problem considerably: the gauge group, dimension, generation count, and discrete symmetries are no longer free variables.

\hypertarget{the-most-constrained-parameters}{%
\subsubsection{3.2 The most constrained parameters}\label{the-most-constrained-parameters}}

Not all 18 parameters are equally constrained. The literature identifies four as critical for complex chemistry:

\textbf{The fine-structure constant α.} Governs atomic binding, molecular bond strengths, and the balance between electromagnetic and thermal energies. Stable atoms require \(\alpha \lesssim 1\) (above this, relativistic effects destabilize inner shells). Complex chemistry requires bond energies \(\gg kT\) at habitable temperatures, giving \(\alpha \gtrsim 10^{-3}\) {[}8{]}. The viable range spans roughly two orders of magnitude: \(10^{-3} \lesssim \alpha \lesssim 10^{-1}\), with our universe at \(\alpha \approx 1/137 \approx 7.3 \times 10^{-3}\) --- within the viable range but not near its center.

\textbf{The electron-to-proton mass ratio \(m_e / m_p\).} Determines the separation of nuclear and atomic scales. Stable atoms with a rich periodic table require \(m_e / m_p \ll 1\) (so electrons orbit far from the nucleus). Complex chemistry breaks down for \(m_e / m_p \gtrsim 10^{-1}\) (electron clouds too compact for directional bonding). The viable range is roughly \(10^{-5} \lesssim m_e / m_p \lesssim 10^{-1}\), spanning four orders of magnitude. Our value: \(m_e / m_p \approx 5.4 \times 10^{-4}\).

\textbf{The quark mass ratio \(m_d - m_u\).} Controls the proton-neutron mass difference and hence nuclear stability. If \(m_d - m_u\) changes sign (proton heavier than neutron), hydrogen is unstable --- catastrophic for chemistry. If \(m_d - m_u\) is too large, no bound nuclei beyond hydrogen exist. The viable range is approximately \(1 \lesssim (m_d - m_u) / \text{MeV} \lesssim 5\) {[}9{]}, about a factor of 5. Our value: \((m_d - m_u) \approx 2.5\) MeV.

\textbf{The strong coupling \(\alpha_s\).} At the QCD scale (\textasciitilde1 GeV), \(\alpha_s\) determines nuclear binding. If \(\alpha_s\) is \textasciitilde10\% smaller, the deuteron is unbound and big bang nucleosynthesis fails to produce helium --- but stellar nucleosynthesis could still produce carbon {[}8{]}. If \(\alpha_s\) is \textasciitilde30\% larger, di-protons are bound, stars burn through hydrogen catastrophically fast. The viable range is roughly \(0.8 \alpha_s^{\text{obs}} \lesssim \alpha_s \lesssim 1.3 \alpha_s^{\text{obs}}\), about a factor of 1.6 {[}10{]}.

\hypertarget{the-viable-fraction}{%
\subsubsection{3.3 The viable fraction}\label{the-viable-fraction}}

The four critical parameters span roughly:

\begin{longtable}[]{@{}lll@{}}
\toprule\noalign{}
Parameter & Viable range (orders of magnitude) & Log-fraction \\
\midrule\noalign{}
\endhead
\bottomrule\noalign{}
\endlastfoot
α & \textasciitilde2 orders within \textasciitilde3 available & \textasciitilde60\% \\
\(m_e / m_p\) & \textasciitilde4 orders within \textasciitilde6 available & \textasciitilde65\% \\
\(m_d - m_u\) & factor of \textasciitilde5 within \textasciitilde10 available & \textasciitilde50\% \\
\(\alpha_s\) & factor of \textasciitilde1.6 within \textasciitilde2 available & \textasciitilde80\% \\
\end{longtable}

The joint viable fraction, treating the parameters as independent (a conservative assumption --- correlations generally enlarge the viable region), is roughly:

\[f_{\text{viable}} \sim 0.6 \times 0.65 \times 0.5 \times 0.8 \approx 16\%\]

This is a rough estimate, but the order of magnitude is robust: the viable fraction is \(\mathcal{O}(10\%)\), not \(\mathcal{O}(10^{-50})\). The remaining 14 SM parameters (heavy quark masses, lepton masses, mixing angles) are less constrained --- they affect the details of nuclear and atomic physics but not the existence of carbon chemistry.

\hypertarget{comparison-to-fine-tuning-claims}{%
\subsubsection{3.4 Comparison to fine-tuning claims}\label{comparison-to-fine-tuning-claims}}

Traditional fine-tuning arguments quote numbers like \(10^{-120}\) (for the cosmological constant) or \(10^{-50}\) (for the Higgs mass). These numbers arise from varying the \emph{structure} --- asking what happens if the gauge group is different, or if there are different numbers of generations, or if d ≠ 3. The OI framework eliminates all of these variations. What remains is a 16\% viable fraction in the critical parameter subspace --- not fine-tuned by any reasonable standard.

The framework's account of the cosmological constant is separate and independent: the \(10^{122}\) discrepancy is the information compression ratio of the trace-out {[}1, §7{]}, not a fine-tuning problem at all. The Higgs mass hierarchy --- why \(m_H \ll M_{\text{Pl}}\) --- is addressed by the lattice structure: the Higgs is a lattice artifact with mass set by the taste-breaking scale {[}2, §7.11{]}, not by radiative corrections to a continuum field.

\begin{center}\rule{0.5\linewidth}{0.5pt}\end{center}

\hypertarget{the-chirality-chain}{%
\subsection{4. The Chirality Chain}\label{the-chirality-chain}}

\hypertarget{from-the-partition-to-parity-violation}{%
\subsubsection{4.1 From the partition to parity violation}\label{from-the-partition-to-parity-violation}}

The OI partition breaks parity: the trace-out over the hidden sector treats left-handed and right-handed spinor components asymmetrically (\(D_{LL} = 0\), \(D_{RR} \neq 0\)) {[}2, §7.8, Theorem 13{]}. This is structural --- any observer in any (S, φ) sees a chiral weak force. The result: the W and Z bosons couple only to left-handed fermions, and parity is maximally violated in all weak processes.

\hypertarget{molecular-parity-violation}{%
\subsubsection{4.2 Molecular parity violation}\label{molecular-parity-violation}}

The weak neutral current (Z boson exchange between electrons and nuclei) produces a parity-violating potential in atoms and molecules. For chiral molecules --- molecules that are not superimposable on their mirror image --- this creates an energy difference between left-handed (L) and right-handed (D) enantiomers {[}11, 12{]}:

\[\Delta E_{\text{PV}} \sim G_F \alpha Z^5 m_e c^2 \left(\frac{a_0}{R}\right)^3\]

where \(G_F\) is the Fermi constant, \(Z\) the nuclear charge, and \(R\) the bond length. For amino acids: \(\Delta E_{\text{PV}} \sim 10^{-14}\) eV \(\sim 10^{-17} \, kT\) at 300 K {[}11{]}. The sign is universal --- it favors L-amino acids --- and is fixed by the partition: the same partition that makes SU(2) left-handed makes L-amino acids lower in energy.

\hypertarget{the-amplification-argument}{%
\subsubsection{4.3 The amplification argument}\label{the-amplification-argument}}

The energy difference \(10^{-17} \, kT\) is far too small to directly select one handedness. The enantiomeric excess at thermal equilibrium would be \(ee \sim 10^{-17}\). The question is whether amplification mechanisms can promote this bias to complete homochirality.

Three mechanisms are experimentally demonstrated. Autocatalytic amplification (the Soai reaction {[}13{]}) amplifies \(ee \sim 10^{-5}\) to \(> 99\%\). Crystallization (Viedma ripening {[}14{]}) drives racemic crystal mixtures to homochirality from small initial biases. Chiral polymerization with cross-inhibition {[}15{]} amplifies small monomer \(ee\) to polymer homochirality.

The gap between \(10^{-17}\) (PVED) and \(10^{-5}\) (Soai threshold) is 12 orders of magnitude. The resolution: the PVED does not need to provide the full \(ee\). In a small prebiotic pool of \(N\) chiral molecules, random statistical fluctuations produce \(ee \sim 1/\sqrt{N}\). For \(N \sim 10^6\) (a microdroplet), statistical \(ee \sim 10^{-3}\) --- within the Viedma threshold. The PVED's role is not to \emph{produce} the excess but to \emph{bias its direction}: in each independent pool, the fluctuation is random in magnitude, but the PVED tips the direction toward L. Over many pools, L-dominant outcomes outnumber D-dominant outcomes. The magnitude of the PVED is irrelevant --- only its sign matters, and the sign is derived.

\hypertarget{the-structural-conclusion}{%
\subsubsection{4.4 The structural conclusion}\label{the-structural-conclusion}}

If the amplification chain is robust --- and every step is either derived, known physics, or experimentally demonstrated --- then biological homochirality is traced to (S, φ). The handedness of every amino acid in every protein in every organism is a consequence of the partition structure that produces quantum mechanics.

\begin{longtable}[]{@{}lll@{}}
\toprule\noalign{}
Step & Content & Status \\
\midrule\noalign{}
\endhead
\bottomrule\noalign{}
\endlastfoot
1 & Partition breaks P & Theorem {[}2, §7.8{]} \\
2 & SU(2)\_L → electroweak PV & Structural consequence \\
3 & Electroweak PV → PVED & Known physics {[}11, 12{]} \\
4 & Statistical fluctuation + directional bias → homochirality & Demonstrated mechanisms {[}13, 14, 15{]}; sign from Step 1 \\
5 & Biological homochirality & Observed \\
\end{longtable}

\begin{center}\rule{0.5\linewidth}{0.5pt}\end{center}

\hypertarget{implications}{%
\subsection{5. Implications}\label{implications}}

\hypertarget{the-dissolution-of-fine-tuning}{%
\subsubsection{5.1 The dissolution of fine-tuning}\label{the-dissolution-of-fine-tuning}}

The fine-tuning problem assumes two freedoms: the laws could have been different, and the parameters could have been different. The OI framework removes the first --- the structure is derived. §3 shows the second is not as constrained as traditionally claimed: the viable parameter fraction is \(\mathcal{O}(10\%)\), not \(\mathcal{O}(10^{-50})\). The apparent fine-tuning was dominated by structural variation that the framework proves doesn't exist.

\hypertarget{the-dissolution-of-the-anthropic-principle}{%
\subsubsection{5.2 The dissolution of the anthropic principle}\label{the-dissolution-of-the-anthropic-principle}}

The anthropic principle assumes there is an ensemble of possible universes from which observation selects. The framework removes the ensemble: the structure is the only structure consistent with embedded observation (characterization theorem {[}1, §3.3{]}; reconstruction theorem {[}2, §10.3{]}). The preconditions for organic chemistry --- atoms, carbon bonding geometry, the scale hierarchy, water, the thermal window, chiral selection --- are structural (§2). The parameters have a viable fraction of \(\mathcal{O}(10\%)\) (§3). There is nothing for observation to select, because there is nothing that could have been otherwise.

\hypertarget{the-autocatalytic-argument}{%
\subsubsection{5.3 The autocatalytic argument}\label{the-autocatalytic-argument}}

The structural chain derives the preconditions for life. This section argues that the transition from prebiotic chemistry to self-replication is not merely possible but statistically expected, given the derived structure.

\textbf{Combinatorial diversity is structural.} Carbon's sp³ hybridization in d = 3 (§2.3) produces chain-forming chemistry with four bonding directions. The number of distinct molecules constructible from C, H, O, N grows exponentially with chain length: \(\sim 20^N\) for amino acid sequences, \(\sim 4^N\) for nucleotide sequences. For chains of length 50: \(\sim 10^{65}\) possible peptides, \(\sim 10^{30}\) possible RNA sequences. This combinatorial explosion is structural --- it follows from the orbital algebra in d = 3 with four valence electrons. No other element in the derived periodic table produces comparable diversity: silicon has similar valence but weaker bonds (\(\sim 2.3\) eV for Si-Si vs.~\(\sim 3.6\) eV for C-C), insufficient for stable long chains at thermal window temperatures.

\textbf{The autocatalytic threshold.} Kauffman's theorem {[}16{]}: in a network of \(N\) molecule types where each molecule has probability \(p\) of catalyzing any given reaction, autocatalytic sets --- closed networks that collectively catalyze their own production from simple feedstocks --- emerge with probability approaching 1 when \(N \times p\) exceeds a critical threshold of \(\mathcal{O}(1)\). The framework provides \(N \sim 10^{65}\) (structural). The question reduces to whether \(p > 10^{-65}\) --- whether catalysis is not astronomically rare among organic molecules.

\textbf{Catalysis is not rare.} Empirically, short random peptides show catalytic activity at rates suggesting \(p \sim 10^{-9}\) to \(10^{-7}\) per pair {[}17{]}. RNA fragments catalyze their own splicing and polymerization {[}18, 19{]}. Simple amino acids catalyze aldol reactions (the Hajos-Parrish reaction). Even individual metal ions in aqueous solution show catalytic activity for a range of organic reactions. With \(p \sim 10^{-9}\) and \(N \sim 10^{65}\): \(N \times p \sim 10^{56}\), exceeding Kauffman's threshold by 56 orders of magnitude.

\textbf{The role of chirality.} The PVED (§4) selects L-amino acids, halving the combinatorial space and --- critically --- eliminating cross-chiral parasitic reactions. In a racemic mixture, L-monomers can be incorporated into D-chains and vice versa, poisoning both. Homochirality, driven by the PVED's directional bias (§4.3), ensures that the combinatorial exploration proceeds within a self-consistent chemical subspace. This is not a minor refinement --- it may be the difference between a productive autocatalytic network and a parasitized one {[}15{]}.

\textbf{What's structural vs.~parametric.}

\begin{longtable}[]{@{}ll@{}}
\toprule\noalign{}
Element & Status \\
\midrule\noalign{}
\endhead
\bottomrule\noalign{}
\endlastfoot
Combinatorial diversity (\(N \sim 20^L\)) & Structural (carbon in d = 3) \\
Energy-driven exploration & Structural (thermal window, §2.7) \\
Chirality eliminates cross-inhibition & Structural (PVED sign, §4) \\
Catalytic probability \(p\) & Parametric (depends on bond energies, activation barriers) \\
\(p > 1/N\) & Empirically satisfied by \textasciitilde56 orders of magnitude \\
\end{longtable}

The framework derives \(N\) (structural), the thermal window (structural), and the chiral bias (structural). The catalytic probability \(p\) is parametric but empirically enormous relative to the threshold. The autocatalytic emergence of self-replication is therefore not a lucky accident but a statistical consequence of the combinatorial diversity that carbon chemistry in d = 3 necessarily produces.

\textbf{Caveats.} Kauffman's theorem assumes well-mixed conditions and does not account for spatial heterogeneity, parasitic side reactions beyond cross-chirality, or the distinction between catalysis and template-directed replication. These are important for the \emph{specific} form life takes (RNA world vs.~metabolism-first vs.~autocatalytic sets) but not for the \emph{existence} of self-sustaining chemical networks. The theorem guarantees that \emph{some} autocatalytic organization emerges; which one dominates is contingent.

\hypertarget{self-replication-as-read-write-cycling}{%
\subsubsection{5.4 Self-replication as read-write cycling}\label{self-replication-as-read-write-cycling}}

\textbf{The structural parallel.} The framework derives QM from read-write cycling: the observer reads the visible sector and writes correlations into the hidden sector through the coupling \(H_{VB}\); the hidden sector stores those writes and returns them on subsequent reads; the resulting statistics are quantum mechanics {[}1, §10.2 of Fundamental{]}. Template replication is the same structural pattern at the molecular level: a polymer reads its own sequence (through complementary hydrogen bond pairing) and writes a copy (through catalytic polymerization). Self-replication is not a special biological capability --- it is molecular read-write cycling, the chemical instantiation of the same information-theoretic pattern that produces quantum mechanics.

\textbf{Von Neumann's threshold.} Von Neumann's self-reproducing automaton theorem {[}20{]} establishes that in any computational system above a complexity threshold, self-reproducing configurations must exist. The threshold requires three capabilities: (a) a universal constructor --- a subsystem that reads instructions and assembles a product; (b) a copying mechanism --- that duplicates the instructions; (c) a control mechanism --- that sequences construction and copying.

The derived chemistry provides all three:

\emph{(a) Universal construction.} Any catalytic polymer that reads a template sequence and assembles a corresponding product from a monomer alphabet. The autocatalytic argument (§5.3) establishes that catalysis is statistically abundant (\(N \times p \sim 10^{56}\)). Template-directed assembly follows from complementary pairing (structural, §2.6) combined with catalytic activity.

\emph{(b) Instruction copying.} Template replication --- a polymer directing the synthesis of its complement through complementary pairing. Three independently structural capabilities whose intersection is template replication (§5.5, below).

\emph{(c) Control.} The thermal window (§2.7) provides environmental cycling --- temperature fluctuations, wet-dry cycles, concentration gradients --- that drives alternation between construction and copying phases. This is structural: the thermal window guarantees that the environment fluctuates on timescales relevant to chemistry.

Von Neumann's threshold is exceeded by the same combinatorial margin (\(\sim 10^{56}\)) that drives the autocatalytic argument. Self-reproducing molecular configurations are not a lucky accident in the derived chemistry --- they are a mathematical consequence of the system's computational richness exceeding von Neumann's threshold.

\hypertarget{from-template-replication-to-evolution}{%
\subsubsection{5.5 From template replication to evolution}\label{from-template-replication-to-evolution}}

The autocatalytic argument (§5.3) establishes that self-sustaining chemical networks are statistically expected. The remaining question: does autocatalysis generically transition to template-directed replication and hence Darwinian evolution?

\textbf{Template replication requires three capabilities.}

\emph{Linear information-carrying polymers.} Carbon's sp³ hybridization (§2.3) allows chain formation: each carbon bonds to the next while retaining side positions for functional groups. A chain of \(N\) monomers carries \(N\) units of sequence information. The possibility of information-carrying linear polymers is structural --- it follows from four-valent bonding in d = 3.

\emph{Complementary pairing.} Hydrogen bonding (§2.6) is directional in d = 3 --- the donor-acceptor geometry is fixed by the orbital algebra. Nitrogen and oxygen have lone pairs capable of forming hydrogen bonds with specific geometries. Two molecules with complementary hydrogen bond patterns pair selectively. This is the structural basis of Watson-Crick pairing (A-T, G-C). The possibility of selective complementary pairing is structural.

\emph{Catalytic copying.} Template-directed catalysis is a subset of all catalysis --- the template positions monomers through complementary pairing, accelerating production of its complement. §5.3 establishes that catalysis is statistically abundant (\(p \sim 10^{-9}\)). Complementary pairing provides the selectivity. Their intersection --- template replication --- is expected whenever both are present in the same molecular system.

\textbf{The RNA precedent.} RNA is simultaneously a linear information-carrying polymer, a substrate for complementary pairing, and a catalyst (ribozymes {[}18{]}). RNA-catalyzed RNA polymerization has been demonstrated experimentally {[}19{]}: an RNA molecule catalyzes the template-directed copying of another RNA molecule. This dual capability is not a biological accident --- in the framework's view, it is the natural intersection of three structural capabilities (linear polymers, complementary pairing, catalytic activity) in a single molecular family.

\textbf{Heritable variation is automatic.} Template copying with finite fidelity produces errors. Errors in a self-replicating polymer are mutations. In a population of replicating polymers competing for monomers: faster replicators outcompete slower ones (selection), copying errors produce variant sequences (variation), and successful variants propagate (heredity). This is Darwinian evolution --- not as a biological principle imposed from outside, but as the inevitable dynamics of a population of imperfect template replicators.

\textbf{The structural status of evolution.}

\begin{longtable}[]{@{}ll@{}}
\toprule\noalign{}
Element & Status \\
\midrule\noalign{}
\endhead
\bottomrule\noalign{}
\endlastfoot
Linear polymers & Structural (carbon sp³ in d = 3) \\
Complementary pairing & Structural (H-bonding geometry in d = 3) \\
Catalytic activity & Statistically expected (§5.3) \\
Template replication & Intersection of the above three \\
Copying errors & Inevitable (finite-fidelity copying) \\
Selection & Inevitable (competition for finite resources) \\
Darwinian evolution & Consequence of replication + variation + selection \\
\end{longtable}

Each row follows from the preceding rows. No row requires contingent facts about our universe beyond the derived structure and the \textasciitilde16\% viable parameter fraction.

\textbf{Caveats.} The argument establishes that Darwinian evolution is \emph{structurally expected} in any system with carbon chemistry, a thermal window, and a chiral aqueous environment. It does not determine: the specific molecular system that first replicates (RNA, PNA, or something else), the timescale for the transition (millions or billions of years), or the specific pathway from simple replicators to cellular organization. These depend on the specific φ, local geochemistry, and contingent history.

\hypertarget{from-evolution-to-intelligence}{%
\subsubsection{5.6 From evolution to intelligence}\label{from-evolution-to-intelligence}}

\textbf{Information processing is generically fitness-enhancing.} In an environment with the derived physics --- a thermal window (§2.7) guarantees fluctuating conditions (\(kT > 0\)) --- organisms that respond to environmental information outcompete those that do not. A cell that detects a nutrient gradient and moves toward it outcompetes one that moves randomly. This is not a contingent fact about Earth --- it follows from the structure: fluctuating environments reward better internal models of external conditions. Selection generically favors information processing.

\textbf{Neural computation is structurally possible.} Fast electrochemical signaling requires ions (Na⁺, K⁺, Ca²⁺ --- present in the derived periodic table) moving through channels in lipid membranes. Carbon chemistry in d = 3 provides amphiphilic molecules (hydrophilic head + hydrophobic tail from the orbital structure) that self-assemble into bilayer membranes. The scale hierarchy (§2.5) separates neural signaling energies (\textasciitilde meV) from chemical bond energies (\textasciitilde eV), so information processing does not destroy the substrate.

\textbf{General intelligence is structurally unconstrained but not guaranteed.} The derived structure \emph{permits} arbitrarily complex neural circuits (the orbital algebra provides materials, the scale hierarchy provides energetic separation), but whether evolution \emph{reaches} general-purpose reasoning --- abstraction, planning, language --- depends on the specific fitness landscape. On Earth, general intelligence arose once in \textasciitilde4 billion years. Whether this is evidence of rarity or inevitability the framework cannot determine. This is the one step in the chain where the structural argument is weakest: information processing is generically favored, but general intelligence is a much more specific capability.

\hypertarget{from-intelligence-to-ai-and-the-self-referential-limit}{%
\subsubsection{5.7 From intelligence to AI and the self-referential limit}\label{from-intelligence-to-ai-and-the-self-referential-limit}}

\textbf{AI is structurally guaranteed given general intelligence.} If general intelligence exists --- for whatever contingent or structural reason --- then AI follows from the derived structure. The reconstruction theorem {[}2, §10.3{]} implies that a sufficiently sophisticated observer can discover (S, φ). An observer that understands its own substrate can build computational devices exploiting it. The derived periodic table provides semiconductors (silicon: element 14, band gap set by the orbital algebra and parametric bond energies). Derived QM provides transistor physics. The scale hierarchy separates computational energies from bond energies. The chain: general intelligence + derived periodic table + derived QM → semiconductor physics → programmable computation → AI.

\textbf{The self-referential loop.} AI is an embedded observer building another embedded observer within the same (S, φ). The characterization theorem {[}1, §3.3{]} applies to the new observer. It sees the same QM, the same ℏ, the same SM, the same dark sector. Not because it is limited by human design, but because the physics is structural --- determined by the partition, not by the observer's complexity. A superintelligent AI has the same observational access as any other embedded observer: it reads the visible sector and writes to the hidden sector through the same coupling. It cannot determine \(h\). It cannot see past the cosmological horizon. It can build better instruments, but it cannot overcome the partition.

The recursion --- AI building AI building AI --- produces ever more sophisticated embedded observers, each subject to the same structural limits. The Gödel-Turing-OI incompleteness family applies at every level: a formal system cannot prove its own consistency (Gödel), a computer cannot decide its own halting (Turing), an embedded observer cannot determine the hidden state (OI). No level of recursive self-improvement overcomes the partition. The limits are not technological --- they are mathematical.

\textbf{What AI can and cannot do.}

\begin{longtable}[]{@{}ll@{}}
\toprule\noalign{}
Capability & Status \\
\midrule\noalign{}
\endhead
\bottomrule\noalign{}
\endlastfoot
Discover (S, φ) from observations & Structurally possible (reconstruction theorem) \\
Build better models of the visible sector & No structural limit \\
Determine the hidden-sector state \(h\) & Provably impossible (characterization theorem) \\
Observe beyond the cosmological horizon & Provably impossible (causal partition) \\
Build faster computers & Structurally possible (derived physics permits it) \\
Overcome the P-indivisibility of QM & Provably impossible (structural, not technological) \\
Recursive self-improvement & Possible within the structural limits \\
\end{longtable}

The framework's prediction is precise: there is no ceiling on the complexity of the observer, but there is a permanent floor on what any observer can access. The ceiling is contingent --- it depends on the specific φ and the specific evolutionary and technological history. The floor is structural --- it is the same for all embedded observers, biological or artificial, at any level of sophistication.

\hypertarget{the-complete-chain}{%
\subsubsection{5.8 The complete chain}\label{the-complete-chain}}

\[(S, \varphi) \xrightarrow{\text{theorem}} \text{QM, d = 3, SM structure}\] \[\xrightarrow{\text{structural}} \text{orbitals, periodic table, carbon, scale hierarchy, water, thermal window, chirality}\] \[\xrightarrow{\text{statistical}} \text{autocatalytic networks (}N \times p \sim 10^{56}\text{)}\] \[\xrightarrow{\text{von Neumann}} \text{self-reproducing configurations (threshold exceeded by same margin)}\] \[\xrightarrow{\text{structural}} \text{template replication (three structural capabilities intersect)}\] \[\xrightarrow{\text{inevitable}} \text{Darwinian evolution (imperfect replication + competition)}\] \[\xrightarrow{\text{generic selection}} \text{information processing} \xrightarrow{\text{structurally possible}} \text{neural computation}\] \[\xrightarrow{\text{contingent}} \text{general intelligence} \xrightarrow{\text{structural}} \text{AI}\] \[\xrightarrow{\text{self-referential}} \text{embedded observers studying embedded observation}\]

One step is contingent: general intelligence. Everything upstream is structural, statistical, or inevitable. Everything downstream is structural given that step.

\hypertarget{what-remains-genuinely-open}{%
\subsubsection{5.9 What remains genuinely open}\label{what-remains-genuinely-open}}

Two questions:

\textbf{Does evolution generically produce general intelligence?} The structural chain makes information processing generically favored, neural computation structurally possible, and the materials for complex circuits available. But general intelligence --- reasoning about reasoning, modeling oneself, asking ``why?'' --- is a specific capability that may or may not be a generic attractor of evolution. On Earth it happened once. The framework cannot determine whether this is typical.

\textbf{What are the ultimate limits of embedded intelligence?} The characterization theorem sets a floor: no observer can determine \(h\), see beyond the horizon, or overcome P-indivisibility. Within these limits, is there a structural bound on the \emph{complexity} of achievable computation? The framework does not currently constrain this --- it says what observers cannot access, not how much they can do with what they can access. Whether there exists a structural bound on computational complexity for embedded observers, analogous to the halting problem for Turing machines, is an open question the framework identifies but does not resolve.

\begin{center}\rule{0.5\linewidth}{0.5pt}\end{center}

\hypertarget{the-existence-question}{%
\subsection{6. The Existence Question}\label{the-existence-question}}

\hypertarget{the-last-question}{%
\subsubsection{6.1 The last question}\label{the-last-question}}

The structural chain derives everything from (S, φ) --- physics, chemistry, life, intelligence, AI. But it does not derive (S, φ) itself. The reconstruction theorem {[}2, §10.3{]} establishes:

\[\text{Observed physics} \quad \longleftrightarrow \quad [(S, \varphi)] / \sim\]

This is conditional: \emph{if} observation exists, \emph{then} (S, φ) exists. The question is whether the framework can say anything unconditional about \emph{why} (S, φ) exists at all.

\hypertarget{mathematical-existence}{%
\subsubsection{6.2 Mathematical existence}\label{mathematical-existence}}

(S, φ) is a finite set and a bijection. This is a mathematical object --- a structure in the space of logical possibilities. It requires no physical substrate, no creator, and no cause. It exists in the same sense that the integers exist or that the permutation group \(S_n\) exists: as a logically consistent structure that cannot fail to be well-defined.

If mathematical structures are taken as necessarily existent --- the Platonic position --- then (S, φ) \emph{necessarily} exists. Not ``happens to exist'' or ``was created'' but cannot fail to exist, because there is no logical contradiction in its definition. On this reading, ``why is there something rather than nothing?'' dissolves: the question assumes non-existence is the default and existence requires explanation. For mathematical structures, existence is the default. Non-existence would require a logical inconsistency, and (S, φ) has none.

\hypertarget{observable-mathematics}{%
\subsubsection{6.3 Observable mathematics}\label{observable-mathematics}}

The framework makes this sharper than generic Platonism. Tegmark's Mathematical Universe Hypothesis {[}7{]} asserts that all consistent mathematical structures are equally real. The OI framework says something more specific: among all mathematical structures, only those containing embedded observers are \emph{observable from within}. The characterization theorem {[}1, §3.3{]} identifies which structures contain embedded observers: exactly those of the form (S, φ) with conditions C1--C3. The reconstruction theorem {[}2, §10.3{]} identifies what those observers see: exactly our physics.

There is no ensemble of equally real structures from which observation selects. There is one equivalence class --- the one compatible with embedded observation --- and it is the one we see. The framework does not need a multiverse, an ensemble, or a selection mechanism. It needs only the observation that (S, φ) is a well-defined mathematical structure and that its observable content is unique.

\hypertarget{the-self-referential-closure}{%
\subsubsection{6.4 The self-referential closure}\label{the-self-referential-closure}}

The complete chain:

\[(S, \varphi) \text{ exists (mathematical structure)}\] \[\xrightarrow{\text{theorem}} \text{contains embedded observers}\] \[\xrightarrow{\text{theorem}} \text{observers see QM, GR, SM}\] \[\xrightarrow{\text{structural}} \text{chemistry, water, thermal window, chirality}\] \[\xrightarrow{\text{statistical}} \text{autocatalytic networks, template replication}\] \[\xrightarrow{\text{inevitable}} \text{evolution} \xrightarrow{\text{generic}} \text{information processing}\] \[\xrightarrow{\text{contingent}} \text{general intelligence} \xrightarrow{\text{structural}} \text{AI}\] \[\xrightarrow{} \text{observers ask: ``why does } (S, \varphi) \text{ exist?''}\]

The answer the framework gives: (S, φ) is a mathematical structure. Mathematical structures exist necessarily. The question can only be asked from inside one. The fact that you are asking it is a structural consequence of the fact that (S, φ) is the kind of structure that contains observers who ask it.

This is not circular --- each step is independently proved or justified. But it is \emph{closed}: the chain returns to its starting point. The framework explains its own explicability. An embedded observer, by the very fact of being embedded, is inside a structure that guarantees the possibility of observers, the possibility of science, and the possibility of the framework itself. The structural chain is not a description of what happened to happen. It is a map of what \emph{must} be observable from within any structure rich enough to contain the question.

\hypertarget{what-cannot-be-answered}{%
\subsubsection{6.5 What cannot be answered}\label{what-cannot-be-answered}}

One question is outside the framework's scope --- not contingently but necessarily:

\textbf{Why is there mathematical structure at all?} The framework takes mathematical existence as given. It does not explain why logical consistency is a property that structures can have, or why ``exists'' and ``is logically consistent'' are related. This is not a gap in the framework --- it is the boundary of all possible frameworks. Any explanation of mathematical existence would itself be a mathematical structure, requiring explanation in turn. The regress has no bottom.

The framework's contribution is not to answer this question but to make it the \emph{only} question. Everything else --- the laws, the constants, the dark sector, the periodic table, life, intelligence, the fine-tuning problem, the anthropic principle --- is derived, dissolved, or identified as generic. The mystery of existence is real. But it is the \emph{only} mystery. And the framework proves it is the \emph{same} mystery for every possible observer in every possible structure.

\begin{center}\rule{0.5\linewidth}{0.5pt}\end{center}

\hypertarget{acknowledgements}{%
\subsection{Acknowledgements}\label{acknowledgements}}

During the preparation of this work, the author used Claude Opus 4.6 (Anthropic) to assist in drafting, refining argumentation, and surveying the relevant literature in nuclear physics, prebiotic chemistry, and origin-of-life research. The author reviewed and edited all content and takes full responsibility for the publication.

\begin{center}\rule{0.5\linewidth}{0.5pt}\end{center}

\hypertarget{references}{%
\subsection{References}\label{references}}

{[}1{]} A. Maybaum, ``The Incompleteness of Observation,'' submitted (2026).

{[}2{]} A. Maybaum, ``The Fundamental Structure of the Observational Incompleteness Framework: From Finite Bijection to the Standard Model,'' submitted (2026).

{[}3{]} P. Ehrenfest, ``In what way does it become manifest in the fundamental laws of physics that space has three dimensions?'' \emph{Proc. Amsterdam Acad.} \textbf{20}, 200 (1917).

{[}4{]} F. R. Tangherlini, ``Schwarzschild field in n dimensions and the dimensionality of space problem,'' \emph{Nuovo Cim.} \textbf{27}, 636 (1963).

{[}5{]} A. D. Sakharov, ``Violation of CP invariance, C asymmetry, and baryon asymmetry of the universe,'' \emph{JETP Lett.} \textbf{5}, 24 (1967).

{[}6{]} C. J. Hogan, ``Why the Universe is just so,'' \emph{Rev.~Mod. Phys.} \textbf{72}, 1149 (2000).

{[}7{]} M. Tegmark, ``Is `the theory of everything' merely the ultimate ensemble theory?'' \emph{Ann. Phys.} \textbf{270}, 1 (1998).

{[}8{]} F. C. Adams, ``The degree of fine-tuning in our universe --- and others,'' \emph{Phys. Rep.} \textbf{807}, 1 (2019).

{[}9{]} A. Jaffe, R. L. Jaffe, and F. Wilczek, ``Quark masses: an environmental impact statement,'' \emph{Phys. Rev.~D} \textbf{79}, 065014 (2009).

{[}10{]} E. Epelbaum, H. Krebs, T. A. Lähde, D. Lee, and U.-G. Meißner, ``Viability of carbon-based life as a function of the light quark mass,'' \emph{Phys. Rev.~Lett.} \textbf{110}, 112502 (2013).

{[}11{]} M. Quack, ``How important is parity violation for molecular and biomolecular chirality?'' \emph{Angew. Chem. Int. Ed.} \textbf{41}, 4618 (2002).

{[}12{]} P. Schwerdtfeger, ``The search for parity violation in chiral molecules,'' in \emph{Relativistic Methods for Chemists} (Springer, 2010).

{[}13{]} K. Soai, T. Shibata, H. Morioka, and K. Choji, ``Asymmetric autocatalysis and amplification of enantiomeric excess of a chiral molecule,'' \emph{Nature} \textbf{378}, 767 (1995).

{[}14{]} C. Viedma, ``Chiral symmetry breaking during crystallization: complete chiral purity induced by nonlinear autocatalysis and recycling,'' \emph{Phys. Rev.~Lett.} \textbf{94}, 065504 (2005).

{[}15{]} G. F. Joyce, G. M. Visser, C. A. A. van Boeckel, J. H. van Boom, L. E. Orgel, and J. van Westrenen, ``Chiral selection in poly(C)-directed synthesis of oligo(G),'' \emph{Nature} \textbf{310}, 602 (1984).

{[}16{]} S. A. Kauffman, \emph{The Origins of Order: Self-Organization and Selection in Evolution} (Oxford University Press, 1993).

{[}17{]} M. R. Ghadiri, J. R. Granja, R. A. Milligan, D. E. McRee, and N. Khazanovich, ``Self-assembling organic nanotubes based on a cyclic peptide architecture,'' \emph{Nature} \textbf{366}, 324 (1993).

{[}18{]} T. R. Cech, ``Self-splicing of group I introns,'' \emph{Annu. Rev.~Biochem.} \textbf{59}, 543 (1990).

{[}19{]} W. K. Johnston, P. J. Unrau, M. S. Lawrence, M. E. Glasner, and D. P. Bartel, ``RNA-catalyzed RNA polymerization: accurate and general RNA-templated primer extension,'' \emph{Science} \textbf{292}, 1319 (2001).

{[}20{]} J. von Neumann, \emph{Theory of Self-Reproducing Automata}, ed.~A. W. Burks (University of Illinois Press, 1966).

\end{document}
